\documentclass[11pt,a4paper]{article}

\usepackage{iftex}
\ifPDFTeX
  \usepackage[T1]{fontenc}
  \usepackage[utf8]{inputenc}
\else
  \usepackage{fontspec}
\fi
\usepackage[spanish,es-noshorthands]{babel}
\usepackage{amsmath,amssymb}
\usepackage{booktabs}
\usepackage{graphicx}
\usepackage{geometry}
\usepackage{float}
\usepackage{caption}
\usepackage{xurl}
\usepackage{hyperref}

\geometry{margin=2.3cm}
\graphicspath{{../results/graficas/}}
\hypersetup{colorlinks=true,linkcolor=black,urlcolor=blue,citecolor=black}
\captionsetup{labelfont=bf,skip=6pt,justification=centering}

% macros generadas; no editar a mano
\newcommand{\Nrep}{14}
\newcommand{\HrmsMean}{0.1809}
\newcommand{\HrmsStd}{0.0056}
\newcommand{\HeqMean}{0.1828}
\newcommand{\EoverN}{-7.938}
\newcommand{\AccMove}{39.2}
\newcommand{\AccArea}{60.7}
\newcommand{\LxMean}{24.92}
\newcommand{\LyMean}{25.90}
\newcommand{\Szero}{2.689}
\newcommand{\NsampTot}{700000}
\newcommand{\RdfPeak}{1.435}
\newcommand{\Qmin}{0.243}
\newcommand{\Qripple}{0.090}
\newcommand{\GqRatio}{0.64}


\title{Resultados y solución de la reimplementación\\[0.35em]
\large Intrinsic ripples in graphene --- $N=240$, $T=300\,\mathrm{K}$}
\author{Emanuel Orozco Gallego}
\date{4 de noviembre de 2026}

\begin{document}
\maketitle

\begin{abstract}
Resultados de \Nrep{} réplicas de Monte Carlo independientes ($N=240$, $T=300\,\mathrm{K}$, $4\times 10^5$ pasos de equilibrado y $5\times 10^5$ de producción). Se comparan con las soluciones del paper: $G_0(q)/N=T/(\kappa S_0 q^2)$ con $\kappa=1.1\,\mathrm{eV}$, el pico de $g(r)$ en el enlace conjugado $1.42\,\text{\AA}$, y la predicción de que una caja de $\sim 25\,\text{\AA}$ no hospeda el máximo de $70\,\text{\AA}$. Las figuras están en PNG.
\end{abstract}

\tableofcontents
\vspace{0.5em}

\section{Protocolo ejecutado}

Se usaron 14 procesos Mojo de 16 núcleos lógicos (2 reservados al sistema, \texttt{taskset} por réplica). Semilla base $42$. Cada réplica escribe CSV en \path{results/datos/}; \path{scripts/generate_plots.py} produce solo PNG/JPG en \path{results/graficas/}.

\begin{center}
\small
\begin{tabular}{lcc}
\toprule
Cantidad & Paper / solución & Esta corrida \\
\midrule
$N$ & 240 (Tabla I) & 240 \\
$T$ & $300\,\mathrm{K}$ (máximo pedido) & $300\,\mathrm{K}$ \\
$L_x\times L_y$ inicial & $24.59\times 25.56\,\text{\AA}$ & $24.59\times 25.56\,\text{\AA}$ \\
Pasos eq.\ /\ prod. & $4\times 10^5$ / $5\times 10^5$ & igual \\
Réplicas & 1 cadena & \Nrep{} independientes \\
Muestras $G(q)$, $g(r)$ & promediado MC & \NsampTot{} \\
\bottomrule
\end{tabular}
\end{center}

\section{Solución 1: ondulaciones térmicas}

El paper (Fig.~1): en equilibrio hay fluctuaciones de altura, no una hoja rígida. Para $N=8640$ citan $h$ típico $\sim 0.6\,\text{\AA}$. En $N=240$ la caja es más pequeña y $\langle h^2\rangle\propto L^2$ en el límite armónico, así que $h_{\mathrm{rms}}$ debe ser menor, pero estrictamente positivo.

\begin{center}
\begin{tabular}{lc}
\toprule
Observable & Valor (media $\pm$ desv.\ entre réplicas) \\
\midrule
$h_{\mathrm{rms}}$ (producción) & $\HrmsMean\pm\HrmsStd\,\text{\AA}$ \\
$h_{\mathrm{rms}}$ al final del eq. & $\HeqMean\,\text{\AA}$ \\
$E/N$ & $\EoverN\,\mathrm{eV/átomo}$ \\
Aceptación desplazamientos & $\AccMove$~\% \\
Aceptación de área & $\AccArea$~\% \\
Caja final $\langle L_x\rangle\times\langle L_y\rangle$ & $\LxMean\times\LyMean\,\text{\AA}$ \\
$S_0$ & $\Szero\,\text{\AA}^2$ \\
\bottomrule
\end{tabular}
\end{center}

$h_{\mathrm{rms}}>0$ confirma rizados intrínsecos. $E/N<0$ indica que la hoja sigue ligada (no se ha fundido ni desarmado a $300\,\mathrm{K}$).

\begin{figure}[H]
\centering
\includegraphics[width=0.82\textwidth]{snapshot_3d.png}
\caption{Instantánea 3D de una réplica. El color es $h_i=z_i-\langle z\rangle$.}
\end{figure}

\begin{figure}[H]
\centering
\includegraphics[width=0.72\textwidth]{snapshot_h_top.png}
\caption{Vista cenital de $h$. Las ondulaciones no están impuestos: salen del MC desde una hoja plana.}
\end{figure}

\begin{figure}[H]
\centering
\includegraphics[width=0.82\textwidth]{traza_equilibrio.png}
\caption{Trazas de $h_{\mathrm{rms}}$ y $E/N$ frente al paso MC. El equilibrado termina en el paso $4\times 10^5$.}
\end{figure}

\begin{figure}[H]
\centering
\includegraphics[width=0.72\textwidth]{histograma_h.png}
\caption{Histograma de $h$ (planos del centro de masa). Distribución ancha y centrada en cero, como describe el paper.}
\end{figure}

\begin{figure}[H]
\centering
\includegraphics[width=0.72\textwidth]{hrms_replicas.png}
\caption{$h_{\mathrm{rms}}$ de cada réplica independiente: la dispersión estima el error de muestreo.}
\end{figure}

\section{Solución 2: $G(q)$ armónico y ausencia del pico a $70\,\text{\AA}$}

La solución analítica de referencia (eqs.\ (5)--(6) del paper) es
\begin{equation}
\frac{G_0(q)}{N}=\frac{T}{\kappa S_0 q^2},\qquad \kappa=1.1\,\mathrm{eV},\quad T=300\,k_B.
\end{equation}
El máximo experimental/numérico del paper está en $q\approx 0.1\,\text{\AA}^{-1}$ ($70\,\text{\AA}$). En esta caja
\[
q_{\min}=\Qmin\,\text{\AA}^{-1} > q_{\mathrm{ripple}}=\Qripple\,\text{\AA}^{-1},
\]
así que \textbf{no hay modos con $q\approx 0.1$}. El paper dice explícitamente que las dos muestras más pequeñas no muestran la caída de $G(q)$ a bajo $q$. Recuperar esa ausencia es la corroboración de la Fig.~3 para $N=240$.

En el intervalo $0.45\lesssim q\lesssim 1.2\,\text{\AA}^{-1}$ (lejos de Bragg y de $q_{\min}$) el cociente mediano $G(q)/G_0(q)$ vale $\GqRatio$. Un valor de orden 1 indica que la convención de Fourier y $\kappa=1.1\,\mathrm{eV}$ describen la hoja a estas escalas.

\begin{figure}[H]
\centering
\includegraphics[width=0.82\textwidth]{Gq_normal_normal.png}
\caption{Correlación normal--normal $G(q)/N$ para $N=240$ (verde) y la recta armónica $G_0(q)$ con $\kappa=1.1\,\mathrm{eV}$ (roja), en la ventana de la Fig.~3 del paper. La muestra sigue $1/q^2$ y se aparta junto a los picos de Bragg $q=2.94$ y $5.11\,\text{\AA}^{-1}$. No aparece el máximo en $q\approx 0.1\,\text{\AA}^{-1}$: esa longitud no cabe en la caja.}
\end{figure}

\section{Solución 3: $g(r)$ y el enlace conjugado}

La Fig.~5 del paper coloca el primer pico de $g(r)$ a $300\,\mathrm{K}$ en $r=1.42\,\text{\AA}$ (enlace conjugado), con hombros hacia $1.31$ y $1.54\,\text{\AA}$ que LCBOPII hace más visibles a $3500\,\mathrm{K}$. Aquí $T\le 300\,\mathrm{K}$, así que el pico debe quedar en el conjugado.

Primer máximo medido: $r=\RdfPeak\,\text{\AA}$ (esperado $1.42\,\text{\AA}$).

\begin{figure}[H]
\centering
\includegraphics[width=0.82\textwidth]{rdf.png}
\caption{$g(r)$ promedio de las réplicas. Verticales: doble, conjugado y simple, como en la Fig.~5 del paper.}
\end{figure}

\section{Pruebas de corroboración}

Las pruebas viven en \path{tests/} y se ejecutan con \texttt{python -m pytest tests}.

\begin{center}
\small
\begin{tabular}{@{}p{0.42\textwidth}p{0.52\textwidth}@{}}
\toprule
Prueba & Qué corrobora \\
\midrule
\path{test_lattice.py} & $N$, $L_x$, $L_y$, coordinación 3, $a=1.42\,\text{\AA}$. \\
\path{test_potential.py} & \texttt{--self-test}: $\Delta E$ local $=$ global; $E/N<0$; Python $=$ C en $N=16$. \\
\path{test_observables.py} & Hoja plana $\Rightarrow G(q)=0$; fórmula $G_0$; pico RDF; $q_{\min}>2\pi/70$. \\
\path{test_mojo.py} & Binario Mojo: $N=240$, $E/N\approx -7.97$ (hoja plana), $\Delta E$ local $=$ global. \\
\path{test_results.py} & CSV: solo $N=240$ y $T=300$; $h_{\mathrm{rms}}>0$; RDF en $1.42$; no hay $q<0.15$; PNG y no PDF. \\
\bottomrule
\end{tabular}
\end{center}

Un modo sintético $h=A\cos(q_x x)$ reproduce $|h_q|$ de la convención del paper (\path{test_observables.py}): la solución~\eqref{eq:G0} no está ``pintada'' a mano, se mide sobre $h_q=\sum_j h_j e^{-iqr_j}$.

\section{Límites de la comparación}

\begin{itemize}
  \item No se simuló $N=8640$ ni $T=3500\,\mathrm{K}$: el pedido fija $N=240$ y $T\le 300\,\mathrm{K}$.
  \item El potencial no es LCBOPII. A esta escala espacial el paper ya predice que el pico de $70\,\text{\AA}$ no aparece; el Tersoff basta para rizados, $G_0(q)$ y $g(r)$ conjugado.
  \item $h_{\mathrm{rms}}$ no debe copiar el $0.6\,\text{\AA}$ de $N=8640$.
\end{itemize}

\section{Archivos de la solución}

\begin{center}
\small
\begin{tabular}{@{}ll@{}}
\toprule
PNG & Contenido \\
\midrule
\path{snapshot_3d.png} & instantánea 3D \\
\path{snapshot_h_top.png} & mapa de $h$ \\
\path{Gq_normal_normal.png} & Fig.~3 restringida a $N=240$ \\
\path{rdf.png} & Fig.~5 a $300\,\mathrm{K}$ \\
\path{traza_equilibrio.png} & convergencia MC \\
\path{histograma_h.png} & $P(h)$ \\
\path{hrms_replicas.png} & error entre réplicas \\
\bottomrule
\end{tabular}
\end{center}

CSV resumen: \path{results/datos/resumen_global.csv}.

\end{document}
